\ifdefined\standalone 
\documentclass[12pt]{article}
\usepackage{graphicx}
\usepackage{float}
\usepackage[a4paper, margin=1in]{geometry}
\usepackage[utf8]{inputenc}
\usepackage{textcomp}
\usepackage{amsmath}
\usepackage{booktabs}
\usepackage{tikz}
\usepackage{enumitem}
\usetikzlibrary{decorations.pathmorphing,patterns}
\usepackage[hidelinks]{hyperref} % <- hypertext

\begin{document}
\fi
\section{Drying and Mass Conservation}

This verification case evaluates the ability of \textsc{Gpyro} to ensure mass conservation during a coupled drying and pyrolysis process. A zero-dimensional simulation is performed on an initial $1~\mathrm{kg}$ sample of wet solid containing $20\%$ moisture by mass. The sample is subjected to a constant heating rate of $1~\mathrm{K\,min^{-1}}$.The process is modeled using two successive reactions. The first reaction represents drying, in which the wet solid releases water:
\begin{equation}
\mathrm{Solid_{wet}} \;\rightarrow\; \theta_1\,\mathrm{Solid_{dry}} + (1-\theta_1)\,\mathrm{H_2O}
\label{reaction:deshydration}
\end{equation}
followed by a one-step pyrolysis reaction of the dry solid:
\begin{equation}
\mathrm{Solid_{dry}} \;\rightarrow\; \theta_2 \mathrm{Char} +(1-\theta_2) \mathrm{Pyrolysate}
\label{reaction:pyrolysis}
\end{equation}

The kinetic parameters of the reactions are given in Table~\ref{tab:drying_0d_rxns_params}.

The expected analytical outcome consists of $0.20~\mathrm{kg}$ of water, $0.48~\mathrm{kg}$ of char, and the remaining mass converted into pyrolysis gases, thereby conserving the initial $1~\mathrm{kg}$ of material.



\begin{table}[H]
\centering
\caption{Reaction parameters for the drying and pyrolysis processes.}
\label{tab:drying_0d_rxns_params}
\renewcommand{\arraystretch}{1.3}
\begin{tabular}{c c c c c c}
\hline
Reaction &
$Z$ ($\mathrm{s^{-1}}$) &
$E$ ($\mathrm{kJ\,mol^{-1}}$) &
$T_{\mathrm{crit}}$ ($\mathrm{K}$) &
$n$ ($-$) &
$\theta$ ($-$) \\
\hline
Eq.~\eqref{reaction:deshydration} & $8.04\times10^{7}$ & 80  & 373.15 & 1 & 0.8 \\
Eq.~\eqref{reaction:pyrolysis}   & $1.97\times10^{15}$ & 188 & --     & 1 & 0.6 \\
\hline
\end{tabular}
\end{table}

The temporal evolution of the individual mass contributions predicted by \textsc{Gpyro} is compared against the analytical solution in Figure~\ref{fig:mass_balance_figs}, demonstrating correct mass conservation throughout the simulation.

\begin{figure}[H]
\centering
\includegraphics[width=1\linewidth]{Mass_drying.png}
\caption{Comparison of \textsc{Gpyro} predictions with analytical mass balance results for the drying verification case.}
\label{fig:mass_balance_figs}
\end{figure}


\ifdefined\standalone
\end{document}
\fi






